\section*{Chapter \ref{ch:mx:request-for-quote-and-dealer-markets-in-theory} --- Request for Quote and Dealer Markets in Theory}

\begin{solution}{exo:mx:request-for-quote-and-dealer-markets-in-theory:1}
$\mu_{lo}=\lambda_ds/(\lambda_u+\lambda_d+\rho)=0.2\times0.8/102.2=0.0016$.
\end{solution}

\begin{solution}{exo:mx:request-for-quote-and-dealer-markets-in-theory:2}
$C(1)=10$ (no competition, no leakage); $C(2)=10-(\sqrt{29}-2)\times0.564+1=9.09$ basis points.
\end{solution}

\begin{solution}{exo:mx:request-for-quote-and-dealer-markets-in-theory:3}
When it equals the adverse selection, 1.5 basis points: both clients are then quoted the dealer's cost, 2.0.
\end{solution}

\begin{solution}{exo:mx:request-for-quote-and-dealer-markets-in-theory:4}
With $z=0$ the dealer passes the whole surplus to the investor: bid and ask both equal the inter-dealer price. The spread is $z(\Delta V_h-\Delta V_l)$,
and the gap between the types' reservation values shrinks when a low type can sell sooner: it bears the holding cost for less time.
\end{solution}

\begin{solution}{exo:mx:request-for-quote-and-dealer-markets-in-theory:5}
The common part adds to the dispersion of bids but each dealer shades for it, so it adds leakage-free noise without lowering the expected winning
price as much: the benefit per extra dealer falls while the leakage per dealer does not.
\end{solution}

\begin{solution}{exo:mx:request-for-quote-and-dealer-markets-in-theory:6}
A peripheral dealer reaches most sellers only through the core, so the bond passes through more dealers (2.58 hops on average against 1.77) and each
charges a markup. Links to more core dealers, or direct links between peripheral dealers, would shorten the chains.
\end{solution}

\begin{solution}{exo:mx:request-for-quote-and-dealer-markets-in-theory:7}
Three dealers, 7.77 basis points (against two dealers and 9.09 with $\sigma_c=2$).
\end{solution}

\begin{solution}{exo:mx:request-for-quote-and-dealer-markets-in-theory:8}
Each dealer asked and not chosen learns that a seller is about and costs the client leakage; beyond a few dealers the leakage grows faster than the
competition helps (at half a basis point per losing dealer, four is best and twelve costs 2 basis points more).
\end{solution}

\begin{solution}{pb:mx:request-for-quote-and-dealer-markets-in-theory:1}
\textbf{1.} See the definition.
\textbf{2.} High and low types switching at $\lambda_u$ and $\lambda_d$; low owners sell and high non-owners buy through dealers met at rate $\rho$;
dealers keep a share $z$ of the surplus.
\textbf{3.} $\lambda_d\mu_{ho}=(\lambda_u+\rho)\mu_{lo}$ with $\mu_{ho}+\mu_{lo}=s$: $\mu_{lo}=\lambda_ds/(\lambda_u+\lambda_d+\rho)$.
\textbf{4.} See section 1; ask $\Delta V_h$, bid $z\Delta V_l+(1-z)\Delta V_h$.
\textbf{5.} It falls: 210 basis points at 10 meetings a year, 36.7 at 100, 4.0 at 1\,000.
\textbf{6.} Dealers submit one price each without seeing the others', and the best price wins at its own price.
\textbf{7.} The winner is the dealer with the highest estimate of the common value; its expected overestimate is $\sigma_c\E[\max\text{ of }n]$.
\textbf{8.} $C(n)=m-(\sqrt{\sigma_c^2+\sigma_p^2}-\sigma_c)\E[\max]+\ell(n-1)$: the markup, the competition among private values, the leakage.
\textbf{9.} Within 0.001 basis points for five dealers.
\textbf{10.} \emph{Named result}: 16 dealers at 0.1 basis point of leakage per losing dealer, 7 at 0.25, 4 at 0.5, 2 at 1, and a single dealer from 2
basis points, with expected costs of 5.52, 6.92, 8.02, 9.09 and 10.00.
\textbf{11.} At 0.5 of leakage: 6 dealers without common-value error, 4 with $\sigma_c=2$, 3 with $\sigma_c=5$.
\textbf{12.} A typical customer contacts few dealers; wider exposure is mitigated by the winner's curse and relationships.
\textbf{13.} Only if leakage per losing dealer is near a quarter of a basis point; measure it before deciding.
\textbf{14.} A few hubs among many peripheral dealers; bonds flow to the core; central dealers charge up to double the round-trip markups but offer
immediacy.
\textbf{15.} Clients of peripheral dealers go through longer chains (2.58 hops against 1.77) and pay 10.4 basis points against 9.2.
\textbf{16.} See the definition; with the chapter's numbers, 3.5 to 0.5 basis points for the informed client as its information's value rises from 0 to
3, against 2.0 for the uninformed.
\textbf{17.} Prices improve with the number of dealers up to about twenty and worsen beyond.
\textbf{18.} Dealers charge lower spreads to their closest counterparties, more so in turmoil; systemically important dealers charge the periphery and
clients more.
\textbf{19.} Compare the price move after RFQs with many and few dealers, or after RFQs that were not traded, against similar bonds without RFQs.
\textbf{20.} The competition of one more bid against the information one more loser takes away.
\end{solution}

\begin{solution}{iq:mx:request-for-quote-and-dealer-markets-in-theory:1}
Few: two to four, chosen for their axes and past pricing, because each losing dealer learns there is a seller in a bond that trades rarely; more when
leakage is low (the bond is liquid) or the private-value dispersion is high.
\iqlookfor{Competition against leakage.}
\end{solution}

\begin{solution}{iq:mx:request-for-quote-and-dealer-markets-in-theory:2}
Winning means having the most optimistic estimate of the common value; a dealer should bid its estimate less the expected overestimate conditional
on winning, which grows with the number of competitors and the estimates' dispersion.
\iqlookfor{Conditioning on winning; shading grows with n.}
\end{solution}

\begin{solution}{iq:mx:request-for-quote-and-dealer-markets-in-theory:3}
The surplus a dealer bargains over is the low type's cost of holding the asset while it waits; easier contact shortens the wait, shrinking the gap
between the types' reservation values and so the spread.
\iqlookfor{Reservation values and waiting.}
\end{solution}

\begin{solution}{iq:mx:request-for-quote-and-dealer-markets-in-theory:4}
Because what it learns from the flow is worth more in its other trades than the adverse selection on this one: \omterm{def:mx:request-for-quote-and-dealer-markets-in-theory:chasing}{information chasing}.
\iqlookfor{Information value against adverse selection.}
\end{solution}

\begin{solution}{iq:mx:request-for-quote-and-dealer-markets-in-theory:5}
Every RFQ with the dealers asked, their quotes and timestamps, the winner, and the market's prices and trades in the bond and related bonds before and
after; so that price moves after RFQs can be compared by number of dealers.
\iqlookfor{Request-level data; before-and-after prices.}
\end{solution}

\begin{solution}{iq:mx:request-for-quote-and-dealer-markets-in-theory:6}
Chains lengthen and markups rise as peripheral dealers route around the weakened core, and clients who need immediacy pay more; after 2008 chains
lengthened by 20\% (Di Maggio, Kermani and Song).
\iqlookfor{Chains, markups, immediacy.}
\end{solution}
